% ===========================================================================
%  第零章 · 预备知识
% ===========================================================================
\setcounter{section}{-1}
\section{预备知识}

本章列出正文将反复使用的三条预备知识。

\block{0.1 比较大小的基本方法：作差}

比较实数 $a$、$b$ 的大小，最基本的方法是考察差 $a-b$ 的符号：
\[
a-b>0 \iff a>b,\qquad a-b=0 \iff a=b,\qquad a-b<0 \iff a<b.
\]
全文所有比较大小、证明不等式的论证，最终都归结为这一条。作差法的长处是便于通分，下面是一个典型例子。

\begin{proposition}[糖水不等式]
用 $a$ kg 白糖制出 $b$ kg 糖水（$b>a>0$），浓度为 $\dfrac ab$；再加入 $m$ kg 白糖（$m>0$），浓度变为 $\dfrac{a+m}{b+m}$。经验告诉我们糖水加糖后变甜，即
\[
\frac{a+m}{b+m}>\frac ab.
\]
\end{proposition}

\begin{proof}
作差、通分：
\[
\frac{a+m}{b+m}-\frac ab=\frac{b(a+m)-a(b+m)}{b(b+m)}=\frac{m(b-a)}{b(b+m)},
\]
分子、分母的每个因子均为正数，故差为正，不等式成立。（书例 1.1-2）
\end{proof}

\sidenote{传递性：由 $A\ge B$ 与 $B\ge C$ 可得 $A\ge C$。本讲中多次以「链式放缩」的形式出现（题 3）。}

\block{0.2 平方的非负性}

由完全平方公式 $(a\pm b)^2=a^2\pm 2ab+b^2$，立即得到本讲义最基本的事实：
\[
(a-b)^2\ge 0,\quad\text{当且仅当 } a=b \text{ 时取等号}.
\]
全文的不等式都以此为源头。另记一个后面要用的恒等式：
\[
(a-b)^2+4ab=(a+b)^2.
\]

\block{0.3 圆周角定理与射影定理}

以下两条定理在板块二与拓展 A 中直接使用，证明要点附后。

\begin{proposition}[圆周角定理（直径情形）]
设 $AB$ 是圆 $O$ 的直径，$D$ 为圆上异于 $A$、$B$ 的点，则 $\angle ADB=90^\circ$。
\end{proposition}

\begin{proof}[证明要点]
$OA=OD=OB$，故 $\triangle OAD$、$\triangle OBD$ 均为等腰三角形，$\angle ODA=\angle OAD$，$\angle ODB=\angle OBD$。在 $\triangle ABD$ 中，
$\angle ADB=\angle ODA+\angle ODB=\angle OAD+\angle OBD$，代入内角和 $180^\circ$，得 $2\angle ADB=180^\circ$，即 $\angle ADB=90^\circ$。
\end{proof}

\begin{proposition}[射影定理（高的情形）]
直角三角形斜边上的高，是垂足所分斜边两段的比例中项：设高为 $h$，两段长为 $p$、$q$，则
\[
h^2=pq.
\]
\end{proposition}

\begin{proof}[证明要点]
高把原三角形分成两个小直角三角形，由锐角互余知两个小三角形相似，对应边成比例 $\dfrac hp=\dfrac qh$，即得 $h^2=pq$。
\end{proof}

\clearpage
